\documentclass[reprint,prx,amsmath,amssymb,aps,superscriptaddress]{revtex4-2}
\usepackage[colorlinks=true,linkcolor=blue,citecolor=blue,urlcolor=blue]{hyperref}
\begin{document}

\title{A reproducibility and evaluation report on GPT-1900}
\author{Brian Keating}
\author{Akshat Singh}
\affiliation{University of California San Diego}
\date{\today}

\begin{abstract}
GPT-1900 is a 3.29-billion-parameter language model trained only on text written before 1900. Its authors evaluated 24 training checkpoints on 8 physics prompts, graded by a hosted language model, and published a mean of 0.93 out of 5 over 192 answers. We repeated the evaluation. The authors' own repeat runs had matched on 2 of 120 paired answers, because the record kept three fields per answer and omitted the seed and the sample count. With both recorded, generation repeats exactly on one machine, 8 of 8 at the published settings. A fresh sweep of all 192 cells, graded blind beside the authors' own published answers by one grader, lands within 0.01 of them: 1.27 against 1.26. The same published answers under the authors' grader carry 0.93. The grader alone moves the score by +0.33, more than a hint that states the answer (+0.28) and about the size of the training gain. On the two general-relativity tasks the hint is worth +0.89 on the 1911 light-bending consequence and +0.00 on the 1907 free-fall seed; given the 1907 observation alone, 0 of 144 answers supply the equivalence. Asked in plain words what a person feels when the elevator cable breaks, no checkpoint at any point in training says that the person feels no weight (0 of 75), a third say that weight is still felt, and the same checkpoints fail Galilean deductions that sit inside their own corpus. The training corpus, 2,074,160 documents, has none dated after 1899. The model reproduces. The number is a property of the grader. The physics is not there.
\end{abstract}

\maketitle

\section{The question}

A published benchmark score is a measurement. It has an instrument, here a hosted language model used as a grader; a sample, here one answer per checkpoint per prompt; and a record, which is what lets a second laboratory repeat it. GPT-1900 \cite{ref1} published 0.93 and a complete archive of generations, and nobody had checked whether the archive could support a repeat. We asked four things. Can the generations be repeated? Does the score survive a change of grader? How much of the score is the wording of the prompt? Is the corpus what it claims to be?

The two general-relativity prompts carry a further interest. Zahavy \cite{ref2} argues that language models carry out deduction from a stated premise but not the abductive step that produces one, with Einstein's 1907 thought about the falling observer as his case. GPT-1900's \texttt{gr\_free\_fall\_equivalence} asks for that step and \texttt{gr\_main\_elevator\_light} asks for its 1911 consequence, so the benchmark contains a direct test. We added a fifth question: when the falling-elevator thought experiment is put to the model in plain words, with no rubric and no hint, what does it say, and does it do any better on ordinary reasoning that its corpus does contain?

\section{What the published record holds}

We read the archive before running any model. Each answer is stored with three fields: prompt, mode, text. Temperature, top\_k, sample count and token limit are in scope at the line that writes the file and are not written. No seed appears in the evaluation script; the engine default of 42 applied silently. Weights are fetched by a mutable name with no revision pinned. The authors' run and its repeat used 1 and 3 samples per prompt; the engine shares one random generator across every row of a batch, so the two could never have matched. 54 of 89 checkpoints in the registry were never published. No repeat run was ever graded.

\section{Methods}

Upstream is pinned at commit \texttt{5b489c8c} and never edited; everything new lives in a separate tools directory that imports it. Each generation is written as one record of 50 fields: run identity, model repository and revision and weight hash, prompt text and hash and condition, seed and sample index and temperature and top\_k and sample count and token limit, device and GPU and dtype as loaded and library versions, the output and its hash and length and stop reason, and eight judge fields. One sample per engine call, so a record is comparable to any other with the same seed regardless of batch size. A comparison tool refuses to report an agreement rate when any generation setting differs between two runs.

Generation ran on SDSC Expanse, one Tesla V100 (32 GB), torch 2.14, weights held as bfloat16 with the format emulated, since the V100 has no bfloat16 hardware; this is recorded per generation. Settings match the published run: temperature 0.7, top\_k 50, seed 42, one sample.

Grading used the authors' rubric and judge prompt unchanged. No API access to the authors' judge model was available, so answers were graded blind: shuffled with a recorded seed, every provenance field stripped, read in fresh sessions of 24 by a Claude model that saw only prompt, rubric and answer, then mapped back through a sealed key. Grades are labelled provisional; a judge failure is recorded as null, never 0. Two batches of each packet were regraded in independent sessions for a reliability figure. Intervals are 95\% percentile bootstraps of the mean, 2,000 resamples.

Open-ended prompts, which have no rubric, were read rather than graded. Each prompt carries a ladder of verdicts fixed before the run: for the gravity prompts, \emph{wrong} (weight is still felt), \emph{none} (never says what is felt), \emph{weightless} (the 1907 observation), \emph{equivalence} (brings in acceleration on its own, the jump), \emph{not-a-force}, \emph{geometry}; for the controls, \emph{correct}, \emph{partial}, \emph{wrong}, \emph{none}. Readers in fresh sessions saw only the prompt, the expected answer, the ladder and the text, and gave the highest rung reached.

\section{Results}

\subsection{Generation repeats on one machine and not across machines}

Same checkpoint, same 8 prompts, temperature 0, run twice: 8 of 8 byte-identical. At the published settings (0.7, seed 42, one sample), run twice: 8 of 8. With deterministic kernels forced: 8 of 8. Normal against deterministic: 8 of 8. The published 118-of-120 divergence is fully explained by the unrecorded seed and the mismatched sample counts.

Across machines it does not repeat. The same checkpoint at temperature 0 on a CPU in float32 against the V100: 0 of 8 identical. The answers share a prefix of 35 to 142 words, then one token differs and the rest diverges. Each machine repeats itself 8 of 8. The authors' GPU code path cannot run in float32 at all; it asserts bfloat16 rotary tables. Device, dtype and library build belong in the record.

\subsection{The full sweep}

All 24 checkpoints, 8 prompts, one sample, published settings: 192 of 192 generated, 0 errors, 64 minutes. Table~\ref{tab:sweep} sets three series side by side.

\begin{table*}[t]
\caption{The sweep under two graders, n=192 each.}
\label{tab:sweep}
\begin{ruledtabular}
\begin{tabular}{lllll}
Answers & Grader & Mean & 95\% interval & 0 / 1 / 2 / 3 / 4 \\
\hline
our sweep & ours & 1.27 & 1.18 to 1.36 & 13 / 126 / 41 / 12 / 0 \\
authors' published & ours & 1.26 & 1.17 to 1.35 & 13 / 125 / 45 / 9 / 0 \\
authors' published & theirs, as published & 0.93 & 0.84 to 1.03 & 44 / 122 / 22 / 3 / 1 \\
\end{tabular}
\end{ruledtabular}
\end{table*}

Our answers against theirs under one grader: 126 of 192 cells exact, mean difference +0.01. Their answers under our grader against their published grades, same bytes: 119 of 192 exact, 188 within one point, +0.33. Rank correlation of the 24 per-checkpoint means, ours against published: 0.71. Training trend, first eight to last eight checkpoints: 0.97 to 1.52 ours, 0.62 to 1.11 published. Grader self-consistency: 96 answers regraded, 89 exact, 96 within one.

The generation step reproduces at the level of the score. The whole gap between 1.27 and 0.93 is the grader: same rubric, same answers, a different Claude, plus a third of a point, mostly from the zero row (44 zeros against 13).

\subsection{Prompt sensitivity}

Three wordings on three checkpoints, three seeds each, n=72 per wording. Scaffolded is the published prompt with the classical assumptions listed; minimal removes them; leading appends one hint paragraph that states the core idea, generated so that every judge field stays byte-identical.

\begin{table}[t]
\caption{Three wordings, eight prompts, n=72 each.}
\label{tab:wording}
\begin{ruledtabular}
\begin{tabular}{lll}
Wording & Mean & 95\% interval \\
\hline
minimal & 1.22 & 1.12 to 1.32 \\
scaffolded & 1.36 & 1.22 to 1.50 \\
leading & 1.64 & 1.46 to 1.83 \\
\end{tabular}
\end{ruledtabular}
\end{table}

The ordering holds on every checkpoint. All 216 answers open with the same sentence, ``You have touched a matter of considerable importance'', regardless of wording; the published sweep has it in 150 of 192.

\subsection{The two general-relativity tasks at n=72}

Both tasks on all 24 checkpoints, three seeds, three wordings (Table~\ref{tab:gr}).

\begin{table*}[t]
\caption{The 1911 consequence and the 1907 seed, n=72 per wording.}
\label{tab:gr}
\begin{ruledtabular}
\begin{tabular}{llllll}
Task & minimal & scaffolded & leading & hint effect, paired & highest score \\
\hline
1911 light bending & 1.31 & 1.57 & 2.46 & +0.89 (+0.69 to +1.08) & 4 \\
1907 free fall & 0.92 & 0.99 & 0.99 & +0.00 (-0.12 to +0.12) & 2 \\
\end{tabular}
\end{ruledtabular}
\end{table*}

\begin{table*}[t]
\caption{The falling elevator put plainly, read blind. Answer counts; looping and incoherent answers, all from the base model, are omitted.}
\label{tab:elev}
\begin{ruledtabular}
\begin{tabular}{lrrrrrr}
Prompt set & n & wrong & none & weightless & not-a-force & equivalence \\
\hline
his question, 25 checkpoints & 75 & 23 & 48 & 0 & 0 & 0 \\
four phrasings, 4 checkpoints & 80 & 26 & 34 & 0 & 0 & 0 \\
1907 observation with help & 60 & 4 & 38 & 1 & 3 & 0 \\
\end{tabular}
\end{ruledtabular}
\end{table*}

On the elevator the hint lifts 39 of 72 answers to 3 or 4; without it none reaches 3. On free fall no answer in 216 reaches 3 under any wording.

The 1907 observation alone, with or without ``what does this tell us about gravity'', 144 answers, read blind: 0 of 144 bring in acceleration, equivalence or geometry on their own. Every automatic keyword hit (7 of 144) was a false positive on reading. The 1911 thought experiment given with no question, 64 continuations across the base model and three checkpoints: 0 reach geometry; the trained model restates ``acceleration and gravity are equivalent'', which is a line of its own prompt, and stops.

The model can be led to the consequence and not to the seed, even when the seed is handed to it. Given the observation alone it never makes the step.

\subsection{The free-fall question asked directly, with controls}

The earlier runs gave the model the 1907 observation. This one gives it only the situation: ``You are standing in an elevator. The cable breaks and the elevator falls freely down the shaft. What do you feel during the fall? Explain why, and say what this tells us about gravity.'' We asked all 24 checkpoints and the base model, three samples each, and readers who did not know the checkpoint or seed placed each answer on the ladder from Methods (Table~\ref{tab:elev}). Of 75 answers, none says the person feels no weight. 48 never say what is felt. 23 say weight is still felt, and that share grows with training, from 3 of 24 in the first eight checkpoints to 12 of 24 in the last eight. The base model only repeats itself.

Four other phrasings, and three prompts that state the observation and add a hint, on four checkpoints with five samples each, give 140 more answers: 0 reach equivalence, 1 reaches weightless, 30 are wrong. When the occupant releases a stone, the trained checkpoints say it falls to the floor in 14 of 15 samples.

We then checked whether the problem is the abductive step alone. Six control prompts, same checkpoints and samples (Table~\ref{tab:controls}): two Galilean deductions the corpus contains (a stone dropped from a ship's mast; two stones of different weight dropped together), one induction (the pendulum law from four data points), one abduction the corpus records (Uranus wanders; an unseen planet), and two beyond the corpus (the Michelson-Morley null result and the black-body peak, without their modern names). The trained checkpoints get the ship's mast right 0 times in 15 and say ``behind the mast'' 10 times; the two stones right 0 times and ``the heavier first'' 9 times; the pendulum law right once; the unseen planet never. The two prompts beyond the corpus score 0.

\begin{table}[tb]
\caption{Reasoning controls, trained checkpoints, 15 answers per prompt. D deduction, I induction, A abduction; ``in'' means the answer is inside the pre-1900 corpus, ``beyond'' that it is not.}
\label{tab:controls}
\begin{ruledtabular}
\begin{tabular}{llrrrr}
Prompt & type & corr. & part. & wrong & none \\
\hline
stone from a ship's mast & D, in & 0 & 2 & 10 & 3 \\
two stones, different weight & D, in & 0 & 6 & 9 & 0 \\
pendulum rule from data & I, in & 1 & 3 & 10 & 1 \\
Uranus wanders & A, in & 0 & 5 & 10 & 0 \\
Michelson-Morley null & A, beyond & 0 & 0 & 6 & 9 \\
black-body peak & A, beyond & 0 & 2 & 4 & 9 \\
\end{tabular}
\end{ruledtabular}
\end{table}

So the model does not fail only the jump. It fails the ordinary steps its corpus contains, and where it commits it gives the pre-Galilean answer in a physicist's voice: the heavier stone first, the stone behind the mast, weight still felt in free fall. Reinforcement learning gave it the voice and the confidence, not the physics.

\subsection{Contamination}

Model side: 16 sentence stems, 11 about post-1900 physics and 5 pre-1900 controls, two checkpoints, greedy and then sampled three times each. Post-1900 stems: 88 of 88 clean. Controls: most fail; the model cannot state Thomson's electron or the Michelson-Morley result either, so the probe measures weak recall in general. Corpus side: the whole training set, 2,074,160 documents, 56 GB, scanned for 15 terms. 0 documents dated 1900 or later. photon, spacetime, Einstein, photoelectric: 0 each. Planck 705, Michelson 277, both at pre-1900 levels. One document dated 1829 describes a 1908 Nobel Prize: a misdated 1909 report, so zero is a lower bound.

\section{What follows}

Four points. A score from an LLM judge is not reproducible unless the judge's model and version are part of the record; here a change of grader moved the published number by more than any prompt manipulation we tried. A record of a generation must hold the seed, the sample count, the model revision, the device and the dtype, or the generation cannot be reproduced by anyone, including its authors. On the benchmark's own general-relativity tasks the model reproduces an idea it is given and does not produce one it is not given; the 0 of 144 is the first stage of the Artificial Einstein Test failing, measured. And the plain-words elevator, with its controls, says why: the model has not got the observation the jump starts from, nor the ordinary physics around it. A rubric score of 1 on this benchmark is an answer that restates the given in the right voice.

\section{Limitations}

Grading. Scores are provisional: a blinded Claude session applying the authors' rubric, not the authors' judge model, which was not reachable without a key. The cross-grade in Sec.~IV.B bounds the effect of that choice but does not remove it. No physicist has graded any answer; the blinded packet exists for that purpose. The open-ended readings are single passes by readers whose counts agree in shape, without a regrade. Regrade agreement on the rubric packets fell from 44 and 45 of 48 exact on the early packets to 35 and 38 of 48 on the later ones.

Hardware. All generation ran on SDSC Expanse on Tesla V100 GPUs, which have no native bfloat16 arithmetic; the published weights are bfloat16 and PyTorch emulates the format. We measured what this costs. The same checkpoint, prompt and seed at temperature 0 on the V100 and on the CPU float32 path gave 0 of 8 identical answers: the texts agree for between 34 and 142 words, a mean of 450 characters, then one near-tied token resolves differently under the two roundings and every later token differs. Each path is exactly repeatable on its own, 8 of 8 on repeat. The authors' GPU code cannot run in float32, since it asserts bfloat16 rotary tables, so a native-bfloat16 run on an H100 is a third numerical path that this allocation could not test, and the authors' own hardware is unrecorded. Our generations therefore reproduce the published score and not the published text, and no reproduction that does not fix device, dtype and library build should expect byte-level agreement.

Coverage. One model and eight rubric prompts, three seeds per cell. The directly posed free-fall question used one phrasing across 24 checkpoints and four phrasings across 4; the controls used 4 checkpoints. 54 of the 89 checkpoints in the authors' registry are unpublished and their results cannot be checked by anyone. The corpus search is lexical and cannot see paraphrase.

\section{Data and code availability}

Every generation, grade and reading is archived under \texttt{evidence/} with its 50-field record; every command's output is under \texttt{logs/}; the lab notebook numbers each result E001 to E050. Code: the upstream repository at commit \texttt{5b489c8c} \cite{ref1}; this study's tools and records will be deposited with a DOI on acceptance.

\section{Acknowledgements}

Computation on SDSC Expanse through ACCESS allocation CIS261668.

\begin{thebibliography}{99}
\bibitem{ref1} M. Hla, GPT-1900, github.com/michaelhla/gpt1900, commit \texttt{5b489c8c} (2026).
\bibitem{ref2} T. Zahavy, Position: LLMs can't jump, ICML 2026 position paper, OpenReview (2026).
\bibitem{ref3} Luo et al., Pretraining language models on historical text (TypewriterLM), arXiv:2606.02991 (2026).
\bibitem{ref4} G. Starace et al., PaperBench: Evaluating AI's ability to replicate AI research, arXiv:2504.01848 (2025).
\bibitem{ref5} Z. S. Siegel et al., CORE-Bench, arXiv:2409.11363 (2024).
\bibitem{ref6} mhla/gpt1900-d34-contradiction-rl-v11, Hugging Face, revision \texttt{71bf395f06dff39ad5ee9f8981b941b3c37daae9}.
\end{thebibliography}
\end{document}
